% Original nested sampling diagram; schematic boxes are not sample counts.
% Final width 112 mm. Task split precedes episode and fragment sampling.
\begin{tikzpicture}[x=\linewidth,y=1mm,
 font=\figurefont\fontsize{8.5}{11}\selectfont,
 task/.style={draw=BookRule,line width=.8pt,fill=BookPaper},
 episode/.style={draw=BookRule,line width=.7pt,fill=white}]
 \foreach \x/\name in {.01/训练任务,.35/验证任务,.69/测试任务} {
  \draw[task] (\x,0) rectangle ++(.30,-57);
  \node[anchor=north] at (\x+.15,-2) {\name};
  \foreach \y/\episode in {-13/回合1,-34/回合2} {
   \draw[episode] (\x+.025,\y) rectangle ++(.25,-18);
   \node[anchor=north] at (\x+.15,\y-1) {\episode};
   \foreach \dx/\seg in {.04/片段,.16/片段} {
    \draw[draw=BookTeal,line width=.7pt,fill=BookTeal!10]
      (\x+\dx,\y-10) rectangle ++(.10,-6);
    \node at (\x+\dx+.05,\y-13) {\seg};
   }
  }
 }
 \draw[BookAmber,dashed,line width=.9pt] (.33,4) -- (.33,-60);
 \draw[BookAmber,dashed,line width=.9pt] (.67,4) -- (.67,-60);
 \node[anchor=south,align=center] at (.50,7)
   {在任务层划分：新任务、组合或因素范围};
 \node[anchor=north,align=center,text=BookMuted] at (.50,-63)
   {同回合片段共享轨迹；同任务回合共享任务参数\\
    片段数量不等于独立测试任务数量};
\end{tikzpicture}
